\documentclass[12pt]{article}
\usepackage[a4paper, margin=1in]{geometry}
\usepackage{mathpazo}
\geometry{left=2.5cm,right=1.5cm,top=2.5cm,bottom=2.5cm}  %margings

\renewcommand{\familydefault}{\sfdefault}
\newcommand{\HI}{\textrm{H\textsc{i}}}

\begin{document}
%===================== FRENCH VERSION ==================
\pagestyle{empty}
	\begin{minipage}[t]{0.48\textwidth}
		\textbf{Imp\'etrant:} \\
		RABARISOA Anjia Randia \\
		\textbf{Encadreur:}\\
		Dr. IANJAMASIMANANA Roger
		
	\end{minipage}
	\hfill
	\begin{minipage}[t]{0.48\textwidth}
		\raggedleft
		Universit\'e d' Antananarivo\\
		Sciences et technologies\\
		Physique et application\\
		Astronomie et Astrophysique\\
	\end{minipage}
	\vspace{1cm}
	\begin{center}
		\large \parbox{\linewidth}{\centering\fontsize{15}{18}\selectfont\linespread{1}\selectfont\bfseries La répartition de l'hydrogène atomique (\HI) cinématiquement froid et chaud et son lien avec la formation stellaire dans les galaxies proches}\\[1em]

		R\'esum\'e
	\end{center}
	
	\vspace{1em}
	
Le processus de formation stellaire au sein des galaxies est aliment\'e par le gaz du milieu interstellaire (MIS), qui comprend des phases moléculaire, atomique et ionisée. Cette \'etude examine la relation entre la formation stellaire et les diff\'erentes phases du MIS dans un \'echantillon de 16 galaxies proches, dont 9 naines/de type tardif et 7 spirales.
Nous avons utilis\'e les données CO de \textit{HERACLES}, les données d'hydrog\`ene atomique (\HI) de \textit{THINGS}, les donn\'ees à 24~$\mu$m de \textit{SINGS}, ainsi que les donn\'ees UV de \textit{GALEX NGS}.
Nous avons décomposé les spectres d'émission de l'\HI\ à l'aide de l'algorithme ROHSA, qui prend en compte la cohérence spatiale des spectres voisins, afin de séparer les composantes \HI\ cinématiquement étroites et larges. Nous utilisons ces composantes comme traceurs cinématiques du milieu neutre froid (CNM) et du milieu neutre tiède (WNM), plutôt que comme des mesures directes des phases thermiques.

Nous examinons la répartition spatiale des composantes cinématiquement froide et chaude et constatons qu'elles s'étendent toutes deux jusqu'aux régions les plus externes des disques atomiques observés. La composante étroite représente en médiane environ 45\% de la masse totale d'\HI\ et présente un aspect plus filamentaire que la composante large. Sa fraction n'est pas systématiquement plus élevée dans les bras spiraux stellaires, contrairement aux attentes théoriques selon lesquelles le refroidissement dans les bras pourrait favoriser la formation de gaz froid. De plus, nous ne trouvons pas d'indication claire que la composante étroite de l'\HI\ soit généralement plus concentrée au voisinage de l'émission CO que la composante large, comme on pourrait s'y attendre si la composante étroite était plus étroitement associée au gaz moléculaire.

Nous examinons la relation de Kennicutt--Schmidt pour les différentes composantes gazeuses afin de déterminer lesquelles sont les plus étroitement liées à la formation stellaire. Dans les galaxies naines, la composante cinématiquement froide présente une corrélation globale plus forte avec le taux de formation stellaire que l'\HI\ total ou la composante chaude. Dans les galaxies spirales, la composante chaude est globalement mieux corrélée au taux de formation stellaire que la composante froide. Cette différence pourrait s'expliquer par le fait que la largeur des raies renseigne sur des conditions du gaz différentes dans ces deux types de galaxies. Dans les naines, l'émission étroite trace probablement du gaz plus froid et moins turbulent, associé à des conditions favorables à la formation stellaire. Dans les spirales, les régions de formation stellaire peuvent contenir du gaz froid dont les raies sont élargies par la turbulence ou des mouvements non résolus. Globalement, le gaz moléculaire reste la composante la plus étroitement liée au taux de formation stellaire, comme l'ont montré les études antérieures. Le rôle du gaz atomique cinématiquement froid dans la transition \HI--H$_2$ reste ambigu à notre résolution sub-kiloparsec.

\vspace{2em}
\noindent\textbf{Mots-clés :} Galaxies proches, Milieu interstellaire, Formation stellaire

%=================++++ENG VERSION-------------------------------------------------------------------

\newpage
\begin{minipage}[t]{0.48\textwidth}
	\textbf{Grantee:} \\
	RABARISOA Anjia Randia \\
	\textbf{Advisor:}\\
	Dr. IANJAMASIMANANA Roger
	
\end{minipage}
\hfill
\begin{minipage}[t]{0.48\textwidth}
	\raggedleft
	University of Antananarivo\\
	Sciences and technology\\
	Physics and application\\
	Astronomy and Astrophysics\\
\end{minipage}
\vspace{1cm}
\begin{center}
		\large \parbox{\linewidth}{\centering\fontsize{15}{18}\selectfont\linespread{1}\selectfont\bfseries The Distribution of Kinematically Cold and Warm Atomic Hydrogen (\HI) and Its Connection to Star Formation in Nearby Galaxies}\\[1em]
	Abstract
\end{center}

\vspace{1em}

The process of star formation within galaxies is fueled by gas from the interstellar medium (ISM), which exists in molecular, atomic and ionised phases. This study investigates the relationship between star formation and the different phases of the ISM in a sample of 16 nearby galaxies, composed of 9 late-type/dwarfs and 7 spirals.
We used CO data from \textit{HERACLES}, atomic hydrogen (\HI) data from \textit{THINGS}, 24 $\mu$m data from \textit{SINGS}, and UV data from \textit{GALEX NGS} surveys. We decomposed the \HI\ emission spectra using the ROHSA algorithm, which takes into account the spatial coherence of neighbouring spectra, to separate kinematically narrow and broad \HI\ components. We use these components as kinematic proxies for the cold neutral medium (CNM) and warm neutral medium (WNM), rather than direct measurements of the thermal phases.\\
We examine the spatial distribution of the kinematically cold and warm components and find that both components extend across the outermost observed atomic discs. The narrow component contains a median of about 45\% of the total \HI\ mass and has a more filamentary appearance than the broad component. Its fractional contribution is not consistently higher in stellar spiral arms, contrary to theoretical expectations that cooling in the arms could favour cold gas formation. In addition, we find no clear evidence that narrow \HI\ is generally more concentrated near CO emission than broad \HI, as would be expected if the narrow component were more closely associated with molecular gas.\\
We examine the Kennicutt--Schmidt relation for the different gas components to investigate which components are more closely linked to star formation. In dwarf galaxies, the kinematically cold component shows a stronger overall correlation with the SFR than the total \HI\ or the warm component. In spiral galaxies, the warm component is more globally correlated with the SFR than the cold component. A possible explanation of this discrepancy is that line width traces different gas conditions in the two types of galaxies. In dwarfs, narrow emission likely traces cooler, less turbulent gas associated with conditions favourable for star formation. In spirals, star-forming regions can contain cold gas whose lines are broadened by turbulence or unresolved motions. Overall, molecular gas remains the component that is more closely linked to the SFR, as previously shown in the literature. The role of kinematically cold atomic gas in the \HI-to-H$_2$ transition remains ambiguous at our sub-kpc resolution.

\vspace{2em}

\noindent
\textbf{Keywords:} 
Nearby galaxies, Interstellar medium, Star formation\\



\end{document}
%